表\ref{tab:practice-mt-blocks}に，各経路条件における第1ブロック（第1～100試行）と第4ブロック（第301～400試行）の平均\(MT\)および比較結果を示す．
各条件の10名について，参加者ごとのブロック平均を用いて対応のある両側\(t\)検定を行い，3条件の検定にHolm補正を適用した．
補正後も，すべての条件で第4ブロックの平均\(MT\)は第1ブロックより有意に短かった（補正後の\(p<0.05\)）．
平均短縮量は，ベースライン条件（C1）で\(190.97\,\mathrm{ms}\)，長距離条件（C2）で\(334.25\,\mathrm{ms}\)，幅狭条件（C3）で\(250.25\,\mathrm{ms}\)であった．

\begin{table}[htbp]
  \centering
  \caption{各経路条件における第1ブロックと第4ブロックの平均\(MT\)および
  対応のある両側\(t\)検定の結果（各条件10名）．}
  \label{tab:practice-mt-blocks}
  \small
  \begin{tabular}{lrrrrrr}
    \hline
    & \multicolumn{2}{c}{平均\(MT\)（ms）} & & & \multicolumn{2}{c}{\(p\)値} \\
    \cline{2-3}\cline{6-7}
    経路条件 & 第1ブロック & 第4ブロック & 平均短縮量（ms） & \(t(9)\) & 未補正 & Holm補正後 \\
    \hline
    ベースライン（C1） &  615.23 &  424.26 & 190.97 & 4.01 & 0.0031 & 0.0062 \\
    長距離（C2）       & 1408.65 & 1074.39 & 334.25 & 4.28 & 0.0021 & 0.0062 \\
    幅狭（C3）         & 1495.28 & 1245.03 & 250.25 & 2.70 & 0.0243 & 0.0243 \\
    \hline
  \end{tabular}
\end{table}
